% 05-signatures.tex: cryptographic evidence.
\section{Authenticating truth: signatures, custody and two witnesses}
\label{sec:signatures}

\subsection{What a signature is for}

A digital signature is a seal that only the holder of a private key can make and that anyone with
the matching public key can check. Polaris signs the SHA3-256 digest of a token value under the
issuing agency's ML-DSA key and stores the signature beside the public key that made it. The claim
a signature carries is narrow and permanent: \emph{this key signed these bytes.} It says nothing
about whether the token is still authorized, which is the next section's problem, and nothing about
who holds the key, which is custody's problem. \sImpl

The algorithm is ML-DSA-65 (FIPS 204, security category 3) by default and ML-DSA-87 (category 5)
beside it; ML-DSA-44 is refused everywhere as below the floor. No signed body hardcodes its
algorithm: the custodied key decides it before signing, because \code{algorithm} is itself a signed
field, and the registry advertises the accepted set. SLH-DSA, the hash-based alternative that does
not rest on lattice assumptions, is registered in the algorithm table so that a rotation away from
lattices is a row update, and no SLH-DSA signer is wired; the seed token filed under it can never
be re-signed, and the posture ledger carries the gap. \sImpl{} for ML-DSA; \sFut{} for a hash-based
signer.

\subsection{Custody}

The private key never sits in application code. It sits behind a custody interface with three
drivers: a file (development), PKCS\#11 (a hardware security module or a software token; the key is
generated inside the token and never leaves it), and a cloud key-management service. Continuous
integration exercises the PKCS\#11 driver against a software token that implements ML-DSA, not a
hardware module, and rotates a key inside it: a token signed under the old in-token key still
verifies after the switch while the new key signs. One production profile makes the module the sole
signer: the application refuses to boot if a file key is present or real signing is unavailable.
Which driver, which module, and who holds the key are decisions the readiness ledger assigns to a
deploying organisation. \sImpl{} for the drivers and the drill; \sExt{} for hardware custody, since
no hardware module has been used.

The placeholder deserves a paragraph, because it is where an honest system is most easily
misread. Without the real-signing flag and a lattice library present, the signing module produces a
deterministic SHA3-256 binding labelled \code{DETERMINISTIC-PLACEHOLDER-SHA3-256}, so that a
developer without the library still exercises every structural invariant. It is not a signature and
cannot be mistaken for one, because the label is checked. Since \code{v9.280} a production boot
fails closed unless real signing is actually available, and the placeholder is a named development
profile that warns loudly when used unnamed. The repository's own history records why: at
\code{v9.58} the headline post-quantum claim was, at the data level, a hardcoded string, and the
real signing module was an unused island until it was wired into issuance. \sImpl

\subsection{Two witnesses at issuance, one at use, sampled}

The two-witness principle states that any cryptographic verdict Polaris ships must be checkable by
a second, independent implementation, written separately, differently represented, sharing no code,
and agreeing on the verdict or abstaining explicitly. A verifier that agrees with itself has
established nothing. At issuance the signature just produced is verified by two independent FIPS 204
implementations, a lattice library and an OpenSSL-backed one, and is refused persistence unless both
accept it; a missing witness at issuance is a refusal, not a downgrade. \sImpl

Verification at use runs one witness, because a stored signature is already known to verify under
both, and one witness detects tampering at about ten times the rate: on the reference machine
roughly 7,800 verifications per second per core against roughly 740 two-witness. The fast path
earns its speed by being continuously checked: a mandatory random fraction of successful checks is
replayed through the second witness, a disagreement increments a counter and pages at the highest
severity, and every response names which witness set ran. Two witnesses that are optional would be
one witness with extra documentation. \sImpl

The same principle governs the zero-knowledge verdict (Section~\ref{sec:privacy}): a Python
re-implementation with plain integers modulo a prime, sharing no code with the Rust prover,
re-derives the membership statement and must agree, and it abstains, in the open, on the one axis it
cannot model, the integrity of the proof bytes. What the principle catches is implementation
divergence. What it does not catch is a shared misreading of the specification, and the repository
says so beside every two-witness claim.

\subsection{Canonical signing: the same bytes on both sides}

Every signed artifact in Polaris is a JSON statement with sorted keys and no whitespace, digested
with SHA3-256, then signed. The application builds the statement; the detached verifier
(Section~\ref{sec:verifiers}) rebuilds it from the artifact to check the signature. The two
constructions must be byte-identical or a genuine artifact fails to verify, and a test, the
canonical-equivalence oracle, pins every signed type on both sides. The one bug the oracle is there
to catch was found in review: an artifact whose \code{algorithm} field was set after signing carried
a signed field the signature did not cover. \sImpl

\subsection{Keys have a lifecycle}

Before \code{v9.328} an authority had one current key and every surface said active. The key
register is now an append-only table of events, registered, retired, or compromised from an instant
that may predate the discovery, from which a signed trust list is published. A verifier holding a
trust list decides a key's status at any instant, so a credential under a compromised issuer key is
rejected independently of the issuer's own manifest, which a thief holding the key could forge, and
a document signed and timestamped before the compromise stays valid while one signed after does not.
The two-instance drill records a compromise on one authority and watches the other's decision flip
over the wire, with no cooperation from the compromised party. Algorithm migration is a key event:
register the ML-DSA-87 key, retire the ML-DSA-65 one from an instant. \sImpl

\subsection{What remains classical}

The token signature, the digests, the Merkle commitments, the zero-knowledge proof and the
client-to-edge key exchange for modern clients are post-quantum. The two internal transport hops
inside the deployment, the certificate signatures, and the operator's WebAuthn credentials remain
classical, each gated on a third party (an OpenSSL version in two container images, the public
certificate ecosystem, authenticator hardware) and each stated in the posture ledger with its
threat and its exposure. The relying-party side already offers ML-DSA-65 first for WebAuthn so
that the day an authenticator implements it, enrolment is post-quantum with no change. None of the
lattice library, the proof system or the circuit has had an external cryptographic audit. \sExt

\begin{layercard}{Layer card: cryptographic evidence}
\qa{What it is}{ML-DSA signatures over SHA3-256 digests, stored with their public keys; a custody interface with file, PKCS\#11 and KMS drivers; two independent verifying implementations; a canonical statement format; an append-only key register and a signed trust list.}
\qa{Why it exists}{A claim that leaves the database needs to carry its own proof, and a proof only one program can check is a promise.}
\qa{What it trusts}{FIPS 204 and 202 as standardised; the two libraries not sharing a bug; the custody backend the operator chose.}
\qa{What it refuses}{A verdict from a single implementation at issuance; an algorithm named in code rather than by the key; a signature stored without its key; a placeholder passing as a signature.}
\qa{What it can see}{The bytes it signs, which are digests: the signing path never sees a person.}
\qa{What it cannot see}{Whether the signed claim is still authorized (Section~\ref{sec:authorization}).}
\qa{If it fails}{A forged signature fails both witnesses; a compromised key is retired from an instant by the register and every verifier holding the trust list rejects what it signed afterwards; a library bug that one witness shares with nobody is caught by the sampled disagreement.}
\qa{Checked by}{\code{check\_pqc\_signing\_wired}, \code{check\_pqc\_second\_witness}, \code{check\_verify\_witness\_sampling}, \code{check\_real\_pqc\_default\_boot}, \code{check\_key\_rotation\_drilled}, \code{check\_trust\_lifecycle}, the canonical-equivalence oracle, the attack suite.}
\qa{Now or later}{Now. The hardware module and the external audit are later, and belong to a deploying organisation.}
\qa{Inside and outside}{Inside: the schema, which stores the signature rows. Outside: the independent verifiers that check them without the issuer.}
\end{layercard}

\nextlayer{A signature the issuer's own server checks proves nothing to a stranger, because the
stranger has to trust the server. The next layer takes the verification out of the issuer's hands
entirely.}
